\chapter*{Lampiran}
\addcontentsline{toc}{chapter}{Lampiran}

% ============================================================
% Lampiran A: Rubrik Penilaian
% ============================================================
\section*{Lampiran A: Rubrik Penilaian Tugas Logika}

\begin{table}[h]
\centering
\small
\begin{tabular}{|>{\raggedright\arraybackslash}p{3cm}|>{\raggedright\arraybackslash}p{3.5cm}|>{\raggedright\arraybackslash}p{3.5cm}|>{\raggedright\arraybackslash}p{3.5cm}|}
\hline
\textbf{Kriteria} & \textbf{Sangat Baik} & \textbf{Baik} & \textbf{Perlu Perbaikan} \\
\hline
Translasi ke Logika & Semua pernyataan benar diterjemahkan & Sebagian besar benar & Banyak kesalahan translasi \\
\hline
Tabel Kebenaran & Lengkap dan benar & Ada sedikit kesalahan & Banyak kesalahan \\
\hline
Bukti Formal & Langkah jelas dan valid & Sebagian langkah kurang jelas & Bukti tidak valid \\
\hline
Induksi Matematika & Basis dan langkah induktif lengkap & Ada kekurangan kecil & Struktur bukti salah \\
\hline
Prolog & Fakta, aturan, query benar & Ada sedikit kesalahan sintaks & Tidak berjalan \\
\hline
\end{tabular}
\caption{Rubrik Penilaian Tugas Logika}
\end{table}

% ============================================================
% Lampiran B: Contoh Template Laporan
% ============================================================
\section*{Lampiran B: Contoh Template Laporan Tugas}
\begin{enumerate}
  \item Judul dan Identitas Mahasiswa
  \item Deskripsi Masalah (pernyataan yang dianalisis)
  \item Translasi ke Simbol Logika (jika applicable)
  \item Tabel Kebenaran atau Langkah Bukti
  \item Hasil Eksekusi Prolog (jika applicable)
  \item Refleksi dan Kesimpulan
\end{enumerate}

% ============================================================
% Lampiran C: Glosarium
% ============================================================
\section*{Lampiran C: Glosarium Istilah Logika}
\begin{itemize}
  \item \textbf{Proposisi}: Pernyataan yang bernilai benar atau salah, tetapi tidak keduanya
  \item \textbf{Negasi}: Penyangkalan dari suatu proposisi; not ($\neg$)
  \item \textbf{Konjungsi}: Gabungan dua proposisi dengan ``dan''; and ($\wedge$)
  \item \textbf{Disjungsi}: Gabungan dua proposisi dengan ``atau''; or ($\vee$)
  \item \textbf{Implikasi}: Pernyataan ``jika ... maka ...''; if-then ($\rightarrow$)
  \item \textbf{Bikondisional}: Pernyataan ``jika dan hanya jika''; iff ($\leftrightarrow$)
  \item \textbf{Tautologi}: Proposisi yang selalu bernilai benar
  \item \textbf{Kontradiksi}: Proposisi yang selalu bernilai salah
  \item \textbf{Kontingensi}: Proposisi yang kadang benar dan kadang salah
  \item \textbf{Kuantifier Universal} ($\forall$): ``Untuk semua''
  \item \textbf{Kuantifier Eksistensial} ($\exists$): ``Terdapat'' atau ``Ada''
  \item \textbf{Modus Ponens}: Aturan inferensi: jika $p \rightarrow q$ dan $p$, maka $q$
  \item \textbf{Modus Tollens}: Aturan inferensi: jika $p \rightarrow q$ dan $\neg q$, maka $\neg p$
  \item \textbf{Induksi Matematika}: Metode bukti untuk pernyataan tentang bilangan bulat nonnegatif
\end{itemize}

% ============================================================
% Lampiran D: Referensi Tambahan
% ============================================================
\section*{Lampiran D: Referensi Tambahan}
\begin{itemize}
  \item Open Logic Project: \url{https://openlogicproject.org/}
  \item forall x: Calgary (Formal Logic): \url{https://forallx.openlogicproject.org/}
  \item Learn Prolog Now!: \url{http://lpn.swi-prolog.org/}
  \item SWI-Prolog: \url{https://swi-prolog.org/}
  \item Discrete Mathematics: An Open Introduction: \url{https://discrete.openmathbooks.org/}
\end{itemize}
